\section{The recursive exploration protocol}\label{sec:protocol}

\subsection{Initialization and target generation}

A run begins by recording its operating mode, initial graph revision, admissible sources, resource budget, and stopping criteria. Goal-directed exploration also pins the overall objective and operational milestone. Open-ended exploration instead specifies a policy for generating and selecting mathematical objectives, such as resolving recorded obstructions, connecting compatible arguments, or developing reusable constructions. Each selected objective and its acceptance conditions are recorded before investigation. An ambitious objective can coexist with a modest operational milestone. A milestone must specify a mathematical or infrastructural deliverable, and changes to its acceptance conditions become explicit events.

Retrieval begins with an active objective or an examination of the collection's open obligations under the objective-generation policy. Candidate connections can arise by composing arguments, transferring a construction to a new setting, strengthening a quantitative bound, weakening a hypothesis, or investigating why a recorded counterexample defeats a route. A system can also propose a conjecture before a researcher has specified a subproblem. The proposal should include a candidate reasoning graph with its missing premises, rather than just a plausible sentence.

Semantic similarity serves as a retrieval aid. Mathematical compatibility requires further work: identify the objects, map the hypotheses, track constants and quantifiers, and show that the proposed conclusion would be useful. A transfer proposal is admitted as a candidate until these bridge obligations are discharged. This makes automatically suggested hypotheses inspectable without presenting them as established results.

\begin{table}[htbp]
\centering\small
\begin{tabularx}{\linewidth}{@{}p{2.6cm}YY@{}}
\toprule
Operation & Question generated & Admission obligation \\
\midrule
Composition & Can one result supply a premise of another? & Align objects, scopes, and quantifiers \\
Transfer & Can a mechanism work in a different setting? & Prove the correspondence and missing bridge lemmas \\
Quantitative refinement & Can a bound meet the next theorem's requirements? & Retain every growing parameter and cost \\
Generalization & Which assumption can be weakened? & Check the modified proof and boundary cases \\
Obstruction analysis & What condition avoids a counterexample? & State exactly what the counterexample rules out \\
Structural synthesis & Do several arguments express one reusable mechanism? & Separate analogy from a proved common statement \\
\bottomrule
\end{tabularx}
\caption{Target-generation operations. They organize candidate questions; none guarantees a useful discovery.}
\end{table}

\subsection{Selection, investigation, and feedback}

Targets are selected according to their relevance to the objective, the obligations they could resolve, feasibility, estimated resource cost, and duplication of earlier work. These are scheduling judgments. A high priority is not a probability that a statement is true. A practical scheduler should retain several different approaches and revisit neglected targets when new evidence changes their prospects.

Investigation alternates constructive and adversarial work. Constructive work develops a proof, representation, algorithm, or example. Adversarial work tests endpoint cases, hidden uniformity assumptions, degenerate parameters, counterexamples in the collection, and possible loss of hypotheses. If a proof requires a new lemma, that lemma becomes a child target. Child targets inherit the parent deadline and the conditions needed by the parent argument.

An investigation can terminate with a proof, a conditional result, a counterexample, a finite computation, a failed attempt, or an unresolved question. Each outcome has a different graph effect. A proof supplies a new candidate inference with evidence. A counterexample blocks a specific statement or route. A failed attempt records a method that did not succeed; it does not prove impossibility. An expired budget produces an open obligation, not a negative theorem.

After review, the graph is updated and affected targets are reconsidered. A new lemma may satisfy a premise of an older route. A counterexample may suggest a weaker statement. A quantitative limitation may redirect attention from one resource to another. In the case study, a finite boundary certificate first motivates a uniform pencil bound. A subsequent target asks how to charge contracted modes when transferring certificates to higher tensor orders. Symmetry reduction already supplies balanced multipliers; the transfer argument establishes the degree factor $n/k$. Classical geometry then extends its application to the complete two-term model.

\begin{figure}[htbp]
\centering
\begin{tikzpicture}[>=Latex,font=\small,
 box/.style={draw=blue!50!black,rounded corners=2pt,fill=blue!4,align=center,text width=4.4cm,minimum height=1.55cm,inner sep=7pt},
 arr/.style={->,thick,draw=blue!50!black}]
\node[box] (objective) at (0,0) {Research policy and current graph\\\footnotesize objectives, scopes, open obligations};
\node[box] (target) at (5.7,0) {Formulate an objective\\or target\\\footnotesize candidate route and admission bounds};
\node[box] (investigate) at (5.7,-2.7) {Investigate and test\\\footnotesize prove, compute, challenge; create subtargets};
\node[box] (review) at (0,-2.7) {Review the outcome\\\footnotesize evidence, assumptions, limits};
\node[box] (update) at (0,-5.4) {Append the graph change\\\footnotesize supported result, refutation, or open obligation};
\node[box] (next) at (5.7,-5.4) {Reassess the frontier\\\footnotesize new questions; milestone or budget check};
\draw[arr] (objective)--(target);
\draw[arr] (target)--(investigate);
\draw[arr] (investigate)--(review);
\draw[arr] (review)--(update);
\draw[arr] (update)--(next);
\draw[arr] (next.east)--++(0.5,0)|-([yshift=0.65cm]target.north)--(target.north);
\node[align=center,font=\footnotesize,text width=10.7cm] at (2.85,-7.4)
 {Human researchers and AI systems can contribute at each stage.\\The event history links the state, action, evidence, and resulting revision.};
\end{tikzpicture}
\caption{The proposed RRE cycle. Arrows describe research activity, not mathematical implication. Feedback may generate new targets and, in the open-ended mode, new overall objectives under the declared policy; it does not permit circular proofs. A stopping decision saves the state and remaining obligations.}
\label{fig:cycle}
\end{figure}


\subsection{A controller specification}

The following pseudocode specifies behavior rather than a particular software architecture. ``Review'' includes source and mathematical assessment at a declared level; it is not assumed to be an infallible algorithm.

\begin{quote}\small
\begin{enumerate}[leftmargin=1.8em]
\item Pin the operating mode, graph revision, source versions, review policy, and deadline, together with an objective and milestone or an objective-generation policy.
\item Record the active objectives and build a frontier from open obligations and candidate connections.
\item While time remains and no stopping condition is met:
  \begin{enumerate}[label=\alph*.,leftmargin=1.6em]
  \item Select a target; record its inputs, rationale, and allocation.
  \item Investigate it; create child targets for missing steps and test the proposed route against relevant obstructions.
  \item Save the outcome, argument, computations, and unresolved obligations.
  \item Review the outcome; assign mathematical status and scoped verification records.
  \item Append the graph change and exploration event; validate identifiers, evidence, and dependencies.
  \item Reassess affected routes and generate the next frontier; if the mode permits new overall objectives, record their rationale and acceptance conditions before investigation.
  \item If the milestone appears satisfied, audit its full acceptance criteria.
  \end{enumerate}
\item Save the final state, actual stopping reason, completed work, and open frontier.
\end{enumerate}
\end{quote}

The controller specification accommodates all four RRE modes. A human can select and investigate targets, collaborators can share these activities, or an AI agent can conduct goal-directed exploration within a fixed charter. In the prospective open-ended mode, the agent would also formulate overall research objectives from the collection under the recorded policy. The operating mode and actor should be recorded for each action; a handoff preserves the current graph, open obligations, and remaining budget. Autonomous exploration may produce candidates and internal checks while still requiring human review before a stronger admission status is granted.

The case study illustrates goal-directed autonomous RRE. The author fixed the objective and authorized its acceptance criteria; the active AI research agent developed successive arguments, generated new targets, and used explicit recording utilities to publish graph revisions. Execution took place through an interactive agent session rather than an unattended controller service. Each saved revision has a separate relation to the recorded events and to the mathematical claims it supports.

\subsection{Significance, milestones, and stopping}\label{sec:significance}

We use provisional significance labels to support research management. They are relative to a declared research program and require a written rationale.

\begin{table}[htbp]
\centering\small
\begin{tabularx}{\linewidth}{@{}p{1cm}p{3cm}Y@{}}
\toprule
Level & Description & Evidence needed for an achieved result \\
\midrule
S0 & Connection or recovery & A useful relationship or known result is correctly identified. \\
S1 & Local advance & A stated subproblem is resolved in its required scope. \\
S2 & Reusable advance & A method or result supports a documented class of further arguments. \\
S3 & Major bottleneck closed & A predeclared major obstacle is resolved with the parameters needed downstream. \\
S4 & Objective resolved & A complete argument establishes the primary objective at the stated review level. \\
\bottomrule
\end{tabularx}
\caption{Significance is separate from proof status and literature novelty. A conjecture may have potential S3 significance without being an achieved S3 result.}
\end{table}

Significance is relative to the declared objective. Here the objective was to close the certificate-degree gap for an infinite family beyond flattenings, with a reusable method and matching bounds. The exact theorem in \cref{thm:main} and its transfer lemma meet those criteria. The result therefore closes the specified bottleneck within tensor-decomposition certificates. The designation ``level 3'' in the run charter expresses this research-management objective; the mathematical contribution is stated explicitly rather than inferred from the label.

A run stops when its declared terminal milestone passes the audit, its time or resource ceiling is reached, the user interrupts it, or execution cannot continue. An open-ended run may continue after an individual objective is resolved if its standing policy and remaining budget permit; ``open-ended'' does not mean unbounded execution. The reason for stopping is recorded literally. A ten-hour ceiling permits an earlier stop after acceptance; it does not imply ten hours of work were performed. Final audit time is included in the budget. Repeatedly restarting a subproblem must not silently reset the deadline.

\subsection{Replay and the limits of completeness}

RRE distinguishes three replay modes. \emph{Research playback} reconstructs recorded questions, decisions, evidence, and graph changes. \emph{Computational replay} reruns saved programs with pinned inputs and compares outputs under declared rules. \emph{Exploration rerun} starts another investigation from the original state and may follow a different path. The first two can be reproducible without the third being deterministic.

A useful event includes identifiers, timestamps, actor attribution, input graph revision, source hashes, target, action, concise rationale, outcome, evidence, and graph change. Model and tool versions, settings, random seeds, and returned candidate proposals should be retained when available. An implementation must report missing fields rather than suggest complete reproducibility from a hash chain alone.

Hashes detect changes relative to preserved digests. The record should also state which activities it captures. The case study saves mathematical targets, arguments, computations, decisions, six graph revisions, and their deltas. It records one concurrent event-log fork and its reconciliation without deleting either branch. These records support scientific playback and certificate replay; they do not purport to capture undocumented alternatives or internal model computation.
